\subsection{Visual synthesis of search-budget behavior and claim stability}
\label{sec:visual-synthesis}

Two additional visual summaries consolidate results already reported above without introducing new model fits, partition searches, or inferential tests. Figure~\ref{fig:budget-audit} visualizes the pre-outcome candidate-budget audit together with the exact test-size pairing used in the primary contrast. Figure~\ref{fig:claim-stability} then summarizes the distinct levels at which the conclusions were stable or protocol-sensitive across the primary analysis, the acyclic-scaffold sensitivity, and the dominant-fragment sensitivity.

\begin{figure}[H]
\centering
\includegraphics[width=0.98\textwidth]{figure5_candidate_budget_audit_v3.pdf}
\caption{Candidate-search budget and exact-size pairing. Panels A and B show the frozen single-group regression budget audit for ESOL and FreeSolv, respectively; the selected target gap continued to decrease through 20,000 candidates, so the production budget was capped rather than interpreted as converged. Panel C shows the separately frozen singleton sensitivity budget audit on a relative scale. Panel D confirms that the paired size-matched and target-balanced protocols use exactly the same test-set cardinality for every primary dataset. These quantities were fixed before predictive outcomes were examined.}
\label{fig:budget-audit}
\end{figure}

\begin{figure}[H]
\centering
\includegraphics[width=0.98\textwidth]{figure6_claim_stability_map_v3.pdf}
\caption{Cross-protocol claim-stability map. The primary classification family contained no corrected target-balanced advantage despite predominantly negative point estimates. The six primary regression cells favored target balancing under the single-group acyclic convention, whereas the singleton sensitivity strongly attenuated the ESOL effects and reversed all three FreeSolv point estimates. Under the dominant-fragment sensitivity, seven of nine classification point-estimate directions reversed while all nine corrected sensitivity inferences remained inconclusive. The figure is a visual synthesis of the reported analyses rather than a new statistical family.}
\label{fig:claim-stability}
\end{figure}
